% ============================================================
% 分片：s03a-genus1  对应原文 PDF 第 21--27 页（印刷页 21--27）
% 负责章节：§3 Explicit Example In Genus One + 3.1--3.5（§3 标题起）
% 译者：trans-s02   审校状态：初稿
% ============================================================

% ==== 原文 PDF 第 21 页（印刷页 21）====

\section{亏格一中的显式例子}

\begin{keytakeaway}
\textbf{核心物理图像：}在亏格 $1$、单分歧点的模型里，内窥的几何图像可以被完全显式地算出：B-模型一侧 ${}^{L}G=SO_3$ 的 Hitchin 模空间含 $4$ 个 $A_1$ 轨形奇点，对应自同构群为 $\mathbb{Z}_2$ 的内窥局部系统；A-模型一侧 $G=SL_2$ 的 Hitchin 模空间本身光滑，但在底中三个特殊点 $w=e_i$ 处的 Hitchin 纤维分裂为两条亏格 $0$ 曲线，正是内窥所需的奇异纤维。
\textbf{数学工具：}椭圆曲线 $y^2=f(x)=x^3+ax+b$ 上的分歧 $SL_2$ Higgs 丛、Hitchin 方程在极点附近的局部形式与单值公式、特征多项式 $\operatorname{Tr}\varphi^2$、$\mathbb{Z}_2$ 不变量 $\rho,y_0^2,h_0^2$ 与模空间方程 (3.15)、判别式与 $A_1$ 奇点的形变/消解、2 阶线丛群 $\mathcal{Q}=\mathbb{Z}_2\times\mathbb{Z}_2$ 在 Hitchin 底与纤维上的作用。
\textbf{主要结论：}对每个 $\sigma_0^2\neq 0$，分歧 $SL_2$ Higgs 丛的模空间是 (3.15) 所描述的复曲面，Hitchin 底由参数 $w$ 给出；三个奇异纤维 $w=e_i$ 各分裂为两个分支 $\boldsymbol{F}_{e_i}^{\pm}$，它们在两个二重点处相交；$\mathcal{Q}$ 通过 2 阶点的平移作用于每个纤维，其中 $T_1$ 保持两个分支不变并在二重点处有不动点，$T_2,T_3$ 在长方向上平移、忠实作用并交换两个分支与两个二重点。
\end{keytakeaway}

\subsection{亏格一的 Higgs 丛}

为构造一个能看出内窥的几何类比的显式例子，我们取 $G=SL_2$，并在亏格 $g_C=1$ 的 Riemann 面 $C$ 上工作，带有单个分歧点 $p$。人们或许会以为分歧会带来额外的复杂性，但事实上，亏格 $1$ 加单个分歧点的情形特别简单，在 Hitchin 纤维化的文献中常被考虑。对 $g_C\ge 2$，相应的模空间维数更高，显式计算困难；对 $g_C\le 1$，$C$ 的基本群是阿贝尔的（或平凡的），在没有分歧时会带来与内窥无关的复杂性，我们愿意在此避开。

正如我们在第 5.3 节所解释的，对 $g_C>1$，虽然显式计算困难，谱曲线方法是强有力的替代手段。但我们更愿意从 $g_C=1$ 的情形开始，因为此时一切都可以直接算出。

首先我们描述亏格 $1$ 的 Riemann 面 $C$ 上无分歧的 Higgs 丛。我们从通常的 $SL_2$ Higgs 丛开始，即对 $(E,\varphi)$，其中 $E$ 是行列式平凡的秩 $2$ 丛。若 $E$ 是半稳定的，它必须具有形式
\begin{equation}\tag{3.1}
E=\mathcal{L}\oplus\mathcal{L}^{-1}
\end{equation}
其中 $\mathcal{L}$ 是次数为 $0$ 的复线丛。若 $\mathcal{L}$ 非平凡，则 $\phi$ 在这个基下必须是
\begin{equation}\tag{3.2}
\phi=
\begin{pmatrix}
a & 0\\
0 & -a
\end{pmatrix}
\end{equation}
其中 $a$ 是 $C$ 上通常的全纯微分。$\mathcal{L}$ 的选择由一条与 $C$ 同构的曲线 $C'$（即 $C$ 的雅可比簇）参数化，而全纯微分空间是一维的。于是我们构造了由 $C'\times C$ 参数化的半稳定 Higgs 丛族，其中 $\mathcal{L}$ 的选择给出 $C'$ 中的点，$a$ 的选择给出 $C$ 中的点。（我们隐含地在 $C$ 上取定了一个特定的全纯微分，以把这样的微分空间等同于 $C$。）然而，把 $\mathcal{L}$ 换成 $\mathcal{L}^{-1}$
% [存疑: 原书 p.21 此句「若 $\mathcal{L}$ 非平凡，则 $\phi$ 必须在这个基下是」与紧接的 (3.2) 把 Higgs 场排作直立的 $\phi$，而同一页上一句的 $(E,\varphi)$ 与全书其余各处（含 (3.8)、(3.10)、(3.12) 等）均排作 $\varphi$。疑为原文记号不一致，此处照原样保留。 来源：原书 p.21 英文原文；存疑待核（A）。]
% [Tutor Note: 推导补足。为何「半稳定」迫使 $E=\mathcal{L}\oplus\mathcal{L}^{-1}$？这是 Atiyah 对椭圆曲线上半稳定丛的分类：亏格 $1$ 上次数 $0$ 的半稳定丛必为其稳定部分的直和，而秩 $2$、行列式平凡的稳定部分只能是两个次数 $0$ 的线丛（稳定线丛即 $SL_2$ 丛的约化）。注意 $a$ 是 $C$ 上全纯微分（即 $H^0(C,K)$ 的元素），在 $C$ 上 $H^0(C,K)$ 一维，这正是后文可以只用一个复参数 $a$ 描述 Higgs 场的原因。 来源：原书 p.21 上下文（A）。]

% ==== 原文 PDF 第 22 页（印刷页 22）====

并改变 $a$ 的符号，会得到同一个 Higgs 丛（同构意义下）。这一操作可以理解为一个规范变换
\begin{equation}\tag{3.3}
\begin{pmatrix}
0 & 1\\
-1 & 0
\end{pmatrix}.
\end{equation}
于是我们可以取 $\mathbb{Z}_2$ 的商，得到由 $(C'\times C)/\mathbb{Z}_2$ 参数化的半稳定 Higgs 丛族。这实际上就是 $C$ 上行列式平凡的秩 $2$ 半稳定 Higgs 丛的模空间。

\subsection{分歧}

现在让我们把分歧纳入进来。在 Higgs 丛的语境中，分歧意味着 [Sim1] $\varphi$ 可以在一个指定的点 $p\in C$（或更一般地在若干这样的点）处有极点，并且其极点部分的特征多项式是预先指定的。我们将考虑单极点的情形。此外，在丛 $E$ 于 $p$ 处的纤维 $E_p$ 中，给定一个 $\varphi$-不变的抛物结构，即一个在 $\varphi$ 的极点部分的作用下不变的旗。此外，这个旗还带有抛物权。整个结构可以通过把 Hitchin 方程改造以纳入奇点而统一地描述。$p$ 附近的行为由取值于 $G$ 的紧实形式的极大环面的 Lie 代数 $\mathfrak{t}$ 中的参数 $\alpha,\beta,\gamma$ 决定。设 $z$ 是 $p$ 附近的局部参数，并写 $z=re^{\ii\theta}$。按 [GW] 的记号，$p$ 附近的局部行为为
\begin{equation}\tag{3.4}
\begin{aligned}
A &= \alpha\, d\theta + \ldots\\
\phi &= \beta\,\frac{dr}{r} - \gamma\, d\theta + \ldots .
\end{aligned}
\end{equation}

由于 Higgs 场 $\varphi$ 正是 $\phi$ 的 $(1,0)$ 部分，这些方程立即决定了它的极点行为（差一个规范等价或共轭）。有 $\varphi\sim\sigma\, dz/z$，其中 $\sigma=(\beta+\ii\gamma)/2$。$\varphi$-不变的抛物结构完全由 $\varphi$ 以及一个包含 $\sigma$ 的 Borel 子群的选择决定。若 $\sigma$ 是正则的（我们将一般地这样假设），则包含 $\sigma$ 的 Borel 子群只有有限多个选择，而选择可由对 $\sigma$ 的特征值排序来确定。因此，在研究具有正则 $\sigma$ 的分歧 Higgs 丛时，抛物结构（或 $\alpha$ 的选择）不需要显式指定。
% [Tutor Note: 推导补足。$\sigma=(\beta+\ii\gamma)/2$ 与 $\varphi\sim\sigma\, dz/z$ 来自把实一形式投影到 $(1,0)$ 部分：取 $z=re^{\ii\theta}$、平坦度量 $|dz|^2$，则 $\star(dr/r)=d\theta$、$\star(d\theta)=-dr/r$，$(1,0)$ 投影为 $\tfrac12(1+\ii\star)$，作用于 $\phi=\beta\,dr/r-\gamma\,d\theta$ 得 $\tfrac12(\beta+\ii\gamma)(dr/r+\ii\,d\theta)=\sigma\,dz/z$。这也解释了为何 $\sigma$ 的两个特征值 $\pm\sigma_0$ 会在 (3.11) 中以 $\sigma_0^2$ 的形式进入特征多项式。 来源：原书 p.22 上下文（A）。]

类似地，我们可以确定具有联络 $\mathcal{A}=A+\ii\phi$ 的局部系统绕 $p$ 的单值。它是
\begin{equation}\tag{3.5}
M=\exp(-2\pi(\alpha-\ii\gamma)).
\end{equation}
% [存疑: 按同一页上文 $\mathcal{A}=A+\ii\phi$ 与 (3.4) 的局部形式直接计算，绕 $p$ 一圈给出 $\exp(+2\pi(\alpha-\ii\gamma))$（取 $z=re^{\mathrm{i}\theta}$、逆时针绕行），而 (3.5) 带有负号。此符号应来自单值的朝向约定（例如把单值定义为 $\exp(-\oint\mathcal{A})$ 或取反向绕行），原书未加说明；此处照原样保留。 来源：原书 p.22 英文原文；存疑待核（A）。]

现在让我们特殊化到亏格 $1$、带一个分歧点的曲线 $C$ 的情形。用代数方程显式描述 $C$ 是方便的：
\begin{equation}\tag{3.6}
y^2=f(x)
\end{equation}
其中我们可以取 $f$ 形如
\begin{equation}\tag{3.7}
f(x)=x^3+ax+b
\end{equation}
并把 $p$ 解释为无穷远点。我们假设 $f$ 的根互不相同，从而 $C$ 光滑。
% [Tutor Note: 易错预警。(3.6) 中的 $(x,y)$ 是 $C$ 上的亚纯函数，$x$ 在无穷远点 $p$ 处有 $2$ 阶极点、$y$ 有 $3$ 阶极点，这正是后文「$x\sim z^{-2}$、$y\sim z^{-3}$」的来源。另注意字母 $a$ 被原书复用：在 (3.2) 中 $a$ 是 $C$ 上的全纯微分（即 $H^0(C,K)$ 的元素），而在 (3.7) 中 $a$ 是椭圆曲线方程的常数系数；阅读 (3.9)--(3.12) 时须分清这两处 $a$。 来源：原书 p.22 上下文（A）。]

% ==== 原文 PDF 第 23 页（印刷页 23）====

为显式描述分歧 Higgs 丛 $(E,\varphi)$，我们先像以前一样取 $E=\mathcal{L}\oplus\mathcal{L}^{-1}$，其中 $\mathcal{L}$ 次数为零。我们可以取 $\mathcal{L}=\mathcal{O}(p)\otimes\mathcal{O}(q)^{-1}$，其中 $q$ 是 $C$ 中某个点，对应于 $(x,y)=(x_0,y_0)$。Higgs 场将具有形式
\begin{equation}\tag{3.8}
\varphi=
\begin{pmatrix}
h & k\\
g & -h
\end{pmatrix},
\end{equation}
其中 $h$ 是 $K\otimes\mathcal{O}(p)$ 的截面，$k$ 是 $K\otimes\mathcal{O}(p)\otimes\mathcal{L}^2$ 的截面，$g$ 是 $K\otimes\mathcal{O}(p)\otimes\mathcal{L}^{-2}$ 的截面。当然，对亏格 $1$ 的 $C$，$K$ 由微分 $dx/y$ 平凡化。相应的线丛都具有一维的全纯截面空间，而 $k,g,h$ 的一般形式为
\begin{equation}\tag{3.9}
\begin{aligned}
h &= \frac{dx}{y}h_0\\
k &= \frac{dx}{y}k_0\left((y-y_0)-\frac{f'(x_0)(x-x_0)}{2y_0}\right)\\
g &= \frac{dx}{y}g_0\left((y+y_0)+\frac{f'(x_0)(x-x_0)}{2y_0}\right)\frac{1}{(x-x_0)^2}.
\end{aligned}
\end{equation}
其中 $h_0,k_0,g_0$ 是复常数，而这些公式是这样得到的。$h$ 的公式无需解释。把 $K$ 平凡化之后，$k$ 应当是 $\mathcal{O}(p)\otimes\mathcal{L}^2=\mathcal{O}(p)^3\otimes\mathcal{O}(q)^{-2}$ 的全纯截面，所以我们寻找一个除在无穷远处有三阶极点之外全纯、并且在 $q$ 处有二阶零点的函数。这样的函数是 $(y-y_0)-f'(x_0)(x-x_0)/2y_0$。类似地，$g$ 应当是 $\mathcal{O}(p)\otimes\mathcal{L}^{-2}=\mathcal{O}(p)^{-1}\otimes\mathcal{O}(q)^2$ 的全纯截面，所以我们寻找一个除在 $q$ 处有二阶极点之外全纯、并在无穷远处消失的函数。这样的函数是 $\left((y+y_0)+\frac{f'(x_0)(x-x_0)}{2y_0}\right)/(x-x_0)^2$。
% [Tutor Note: 推导补足。这里隐含的计数是：$\mathcal{L}=\mathcal{O}(p)\otimes\mathcal{O}(q)^{-1}$ 给出 $\mathcal{L}^2=\mathcal{O}(2p)\otimes\mathcal{O}(q)^{-2}$，故 $\mathcal{O}(p)\otimes\mathcal{L}^2=\mathcal{O}(3p)\otimes\mathcal{O}(q)^{-2}$；而在亏格 $1$ 上 $h^0(\mathcal{O}(3p)\otimes\mathcal{O}(q)^{-2})=3-2+1-1=1$（Riemann--Roch 且 $h^1=0$），所以截面唯一到常数倍，这正是 $k_0$ 这个自由常数的由来。$g$ 的情形同理：$\deg(\mathcal{O}(p)\otimes\mathcal{L}^{-2})=-1$，$h^0(-p+2q)= -1+2+1-1=1$。$\mathcal{O}(p)^3$ 是 $\mathcal{O}(3p)$ 的简写。 来源：原书 p.23 上下文（A）。]

下一步是计算 $\varphi$ 的特征多项式。对 $G=SL_2$，这简单地意味着我们应当计算 $\operatorname{Tr}\varphi^2$，而在当前情形它结果为
\begin{equation}\tag{3.10}
\operatorname{Tr}\varphi^2=\left(\frac{dx}{y}\right)^2\left(2h_0^2+2k_0g_0\left(x+2x_0-\frac{f'(x_0)^2}{4f(x_0)}\right)\right).
\end{equation}
$\operatorname{Tr}\varphi^2$ 的极点部分简单地是 $(dx/y)^2 2xk_0g_0$，它在无穷远处有二阶极点。若 $z$ 是无穷远处的局部参数，我们有 $x\sim z^{-2}$、$y\sim z^{-3}$，且 $(dx/y)^2x\sim 4(dz/z)^2$。$\varphi$ 的极点部分应当共轭于 $\sigma dz/z$，所以我们要求 $\operatorname{Tr}\varphi^2\sim\operatorname{Tr}\sigma^2(dz/z)^2=2\sigma_0^2(dz/z)^2$，其中我们把 $\sigma$ 的特征值记作 $\pm\sigma_0$。于是我们令 $k_0g_0=\sigma_0^2/4$，并写出
\begin{equation}\tag{3.11}
\operatorname{Tr}\varphi^2=\left(\frac{dx}{y}\right)^2\left(2h_0^2+\frac{\sigma_0^2}{2}\left(x+2x_0-\frac{f'(x_0)^2}{4f(x_0)}\right)\right).
\end{equation}
% [Tutor Note: 推导补足。(3.10) 中被跳过的一步是一个三次恒等式：对 $f(x)=x^3+ax+b$ 有 $f(x)-f(x_0)-f'(x_0)(x-x_0)=(x-x_0)^2(x+2x_0)$。把 (3.9) 代入 $\operatorname{Tr}\varphi^2=2h^2+2kg$ 后，$kg$ 含因子 $\left[(y-y_0)(y+y_0)-\frac{f'(x_0)^2(x-x_0)^2}{4y_0^2}\right]/(x-x_0)^2$，利用 $y^2=f(x)$、$y_0^2=f(x_0)$ 及上述恒等式即得 (3.10)。再由 $x\sim z^{-2}$、$y\sim z^{-3}$ 得 $\frac{dx}{y}\sim -2\,\frac{dz}{z}$，故 $(dx/y)^2x\sim 4(dz/z)^2$，这正是定出 $k_0g_0=\sigma_0^2/4$ 的依据。 来源：原书 p.23 上下文（A）。]
% [Tutor Note: 易错预警。(3.10)--(3.11) 中的 $f(x_0)$ 出现在分母，因此在 $f(x_0)=0$（即 $y_0=0$）处公式失效；这也正是下文「$y_0=0$ 或 $y_0=\infty$ 时 (3.9) 失效」的代数根源。另注意 $\operatorname{Tr}\varphi^2$ 作为 Hitchin 映射的分量，其归一化依赖 $\sigma$ 的特征值差：$\operatorname{Tr}\sigma^2=2\sigma_0^2$，若把特征值写成 $\pm\sigma_0$ 则系数 $2$ 不可漏掉。 来源：原书 p.23 上下文（A）。]

只有乘积 $k_0g_0$ 被确定，其原因在于丛 $E=\mathcal{L}\oplus\mathcal{L}^{-1}$ 有自同构群 $\mathbb{C}^\times$，它分别在两个直和项上作用为乘以 $\lambda$ 与 $\lambda^{-1}$，并把 $k_0$ 与 $g_0$ 变换为 $k_0\to\lambda^2k_0$、$g_0\to\lambda^{-2}g_0$。

公式 (3.9) 在 $y_0=0$ 时失效（因为 $y_0$ 出现在 $k$ 与 $g$ 的公式的分母中），或者在 $y_0=\infty$ 时失效（因为 $k$ 与 $g$ 的公式还含有关于 $y_0$ 线性的项），等价地，在 $\mathcal{L}$ 为 $2$ 阶时失效。发生这种情况是因为当 $\mathcal{L}$

% ==== 原文 PDF 第 24 页（印刷页 24）====

为 $2$ 阶时，我们有 $\mathcal{L}\cong\mathcal{L}^{-1}$，并且（若 $\sigma\ne 0$）不存在稳定的或半稳定的分歧 Higgs 丛 $(E,\varphi)$ 满足 $E=\mathcal{L}\oplus\mathcal{L}^{-1}=\mathcal{L}\oplus\mathcal{L}$。事实上，取这样的 $E$，$K\otimes\mathrm{ad}(E)\otimes\mathcal{O}(p)\cong K\otimes\mathcal{O}(p)^{\oplus 3}$ 的全纯截面在 $p$ 处不可能有极点，从而条件 $\operatorname{Tr}\varphi^2\sim 2\sigma_0^2(dz/z)^2$ 无法被满足。反之，若 $\mathcal{L}$ 为 $2$ 阶且 $\sigma\ne 0$，则必须取 $E$ 为 $\mathcal{L}$ 被 $\mathcal{L}$ 的非平凡扩张。
% [Tutor Note: 推导补足。「不可能有极点」是次数计算：$\mathcal{L}^2\cong\mathcal{O}$ 时 $\mathrm{ad}(E)\cong\mathcal{O}^{\oplus 3}$，故 $K\otimes\mathrm{ad}(E)\otimes\mathcal{O}(p)\cong K\otimes\mathcal{O}(p)^{\oplus 3}$；而 $\deg(K\otimes\mathcal{O}(p))=2g-2+1=1$，由 Riemann--Roch 得 $h^0=1$、$h^1=0$，其全纯截面（到常数倍）由 $dx$ 给出、在 $p$ 处正则，不可能带有极点，与 $\operatorname{Tr}\varphi^2$ 在 $p$ 处有二阶极点矛盾。这正是为何 $\mathcal{L}$ 为 $2$ 阶时必须改用非平凡扩张。 来源：原书 p.24 上下文（A）。]

\subsection{模空间}

对参数 $\sigma_0^2$ 的每个选择，我们构造了一个分歧 Higgs 丛族。底丛是 $E=\mathcal{L}\oplus\mathcal{L}^{-1}$，而 Higgs 场 $\varphi$ 已如上描述。$E$ 的选择依赖于点 $q$，或等价地依赖于对 $(x_0,y_0)$（满足 $y_0^2=f(x_0)$），而 $\varphi$ 的选择还额外依赖于参数 $h_0$。要在这种情况下构造分歧 Higgs 丛的模空间 $\mathcal{M}_H$，我们必须考虑对换 $\tau:\mathcal{L}\to\mathcal{L}^{-1}$，它作用为 $(x_0,y_0,h_0)\to(x_0,-y_0,-h_0)$。

对 $(q,h_0)$ 或三元组 $(x_0,y_0,h_0)$ 定义了 $C\times C$ 中的一个点，而在容许对称 $\tau$ 之后，模空间似乎是 $(C\times C)/\mathbb{Z}_2$，与 $\sigma_0$ 无关。这实际上并不完全正确，原因在于第 3.2 节末尾所指出的那一点。空间 $(C\times C)/\mathbb{Z}_2$ 有四个 $A_1$ 奇点，位于 $h_0=0$ 且 $q$ 为 $2$ 阶的点处。对 $\sigma\ne 0$，$A_1$ 奇点被形变，模空间变得光滑。我们没有看到这一形变，因为我们的分析没有覆盖 $q$ 为 $2$ 阶的情形。若 $\sigma=0$ 而 $\alpha\ne 0$，则 $A_1$ 奇点是被消解而不是被形变；从超凯勒的观点看，对 $\sigma$ 与 $\alpha$ 的一般取值，这一现象实际上是同时的形变与消解，如 [Kr] 中那样。本文不会直接描述这一形变或消解，但在下文方程 (3.15) 中，借助 Hitchin 纤维化描述了形变后的模空间。
% [Tutor Note: 易错预警。这里区分了两种局部模型：$A_1$ 奇点的**形变**（deformation，$\sigma\ne 0$，对应复结构 $I$ 下一族光滑超凯勒流形）与**消解**（resolution，$\sigma=0$ 而 $\alpha\ne 0$，对应把 $\mathbb{C}^2/\mathbb{Z}_2$ 的奇点用一条 $\mathbb{P}^1$ 替换）。原文说对一般 $(\sigma,\alpha)$ 是「同时的形变与消解」，这正是超凯勒旋转下的图像：$I$ 下看到的形变与 $J$ 下看到的消解由同一个超凯勒度量联系起来（[Kr]）。切勿把二者混为一谈，也别把 (3.15) 误当作消解后的空间——它只描述形变。 来源：原书 p.24 上下文（A）。]

这里或许应当就分歧在亏格 $1$ 中的作用作一点评论。对 $g_C>1$，一般的半稳定丛实际上是稳定的；但对 $g_C=1$（且无分歧）以及单连通的 $G$，不存在严格稳定的丛。最接近的是半稳定丛，例如 $G=SL_2$ 时的 $E=\mathcal{L}\oplus\mathcal{L}^{-1}$。同样地，一般地 $\mathcal{M}_H$ 参数化稳定与半稳定 Higgs 丛。对 $g_C=1$，无分歧的 Higgs 丛充其量是半稳定的，所以 $\mathcal{M}_H$ 实际上参数化半稳定 Higgs 丛。有了分歧，情况就变了。带一个分歧点的半稳定 Higgs 丛实际上是稳定的（若 $\sigma\ne 0$）。其原因在于 $E=\mathcal{L}\oplus\mathcal{L}^{-1}$ 的潜在失稳层（即直和项 $\mathcal{L}$ 与 $\mathcal{L}^{-1}$）不是 $\varphi$-不变的，因此不违反 Higgs 丛 $(E,\varphi)$ 的稳定性。

\subsection{Hitchin 纤维化}

接下来我们需要理解 Hitchin 纤维化。Hitchin 纤维化就是把 Higgs 丛 $(E,\varphi)$ 映到 $\varphi$ 的特征多项式的映射。对 $G=SL_2$，这个特征多项式约化为 $\operatorname{Tr}\varphi^2$。在当前语境中，按 (3.11)，该特征多项式是 $x$ 的一次多项式（乘以 $(dx/y)^2$）。$x$ 的一次项系数是固定的，而 Hitchin 纤维化就是提取常数项的映射，我们把这个常数项记作 $w_0$。换言之，Hitchin 纤维化把 Higgs 丛 $(E,\varphi)$ 映到
\begin{equation}\tag{3.12}
w_0=2h_0^2+\frac{\sigma_0^2}{2}\left(2x_0-\frac{f'(x_0)^2}{4f(x_0)}\right).
\end{equation}

% ==== 原文 PDF 第 25 页（印刷页 25）====

Hitchin 纤维化的纤维由满足 (3.12) 与
\begin{equation}\tag{3.13}
y_0^2=f(x_0)
\end{equation}
的变量 $x_0,y_0,h_0$ 描述，并且还受 $\mathbb{Z}_2$ 对称 $(x_0,y_0,h_0)\to(x_0,-y_0,-h_0)$ 的作用。除 $x_0$ 之外，$\mathbb{Z}_2$ 不变量是 $y_0^2$、$h_0^2$ 以及
\begin{equation}\tag{3.14}
\rho=(2/\sigma_0)y_0h_0.
\end{equation}
（因子 $2/\sigma_0$ 是为方便而引入的。）其中，按最后两个方程，$y_0^2$ 与 $h_0^2$ 等价于 $x_0$ 的有理函数，因而可以略去；而 $\rho$ 满足一个二次方程。若此后我们把 $x_0$ 写作 $u$，并令 $w_0=-w\sigma_0^2/2$，则 $\rho$ 所满足的方程为
\begin{equation}\tag{3.15}
\rho^2=-(2u+w)f(u)+\frac{f'(u)^2}{4}.
\end{equation}
% [Tutor Note: 推导补足。(3.15) 由 (3.12)--(3.14) 消去 $x_0,y_0,h_0$ 得到，关键一步是：由 (3.12) 解出 $h_0^2=\frac{w_0}{2}-\frac{\sigma_0^2}{4}\left(2u-\frac{f'(u)^2}{4f(u)}\right)$，代入 (3.14) 得 $\rho^2=\frac{4}{\sigma_0^2}y_0^2h_0^2=\frac{4}{\sigma_0^2}f(u)h_0^2$；再用 $w_0=-w\sigma_0^2/2$ 化简，$\sigma_0^2$ 恰好约去，即得 (3.15)。这也解释了原书为何强调 $\sigma_0$ 不出现：$\sigma_0$ 可通过重新标度 $\varphi$ 消去，(3.15) 只依赖 $w$ 与 $u$。 来源：原书 p.25 上下文（A）。]

关于复变量 $\rho,u,w$ 的这个方程描述了一个复曲面，它就是分歧 $SL_2$ Higgs 丛的模空间 $\mathcal{M}_H$。（在这种写法中，$u,\rho=\infty$ 处的某些点被略去了。）有一个简单的解释说明它为何不依赖于刻画分歧的参数 $\alpha,\beta,\gamma$。复结构 $I$ 不依赖 $\alpha$，而只要只有一个分歧点，$\sigma=(\beta+\ii\gamma)/2$ 就可以通过重新标度 $\varphi$ 消去（只要它不为零），正如我们所做的那样。

若把 $w$ 取为固定的复数，我们得到一条代数曲线 $\boldsymbol{F}_w$，它是 Hitchin 纤维化的纤维。右端是关于 $u$ 的四次多项式，并且对一般的 $w$，$\boldsymbol{F}_w$ 是亏格 $1$ 的光滑曲线。

$\boldsymbol{F}_w$ 何时奇异？这恰好发生在多项式 $g(u)=(2u+w)f(u)-f'(u)^2/4$ 的两个根重合时，换言之，在其判别式为零时。设 $e_1,e_2,e_3$ 是多项式 $f(u)=u^3+au+b$ 的根。（当然 $e_1+e_2+e_3=0$，并且我们假设 $e_i$ 互不相同，从而 $C$ 光滑。）$g$ 的判别式为 $(e_1-e_2)^2(e_2-e_3)^2(e_3-e_1)^2(w-e_1)^2(w-e_2)^2(w-e_3)^2$，所以 $\boldsymbol{F}_w$ 奇异恰好当 $w$ 等于某个 $e_i$。例如，若 $w=e_1$，我们发现
\begin{equation}\tag{3.16}
-(2u+w)f(u)+\frac{f'(u)^2}{4}=\frac{1}{4}\left((u-e_1)^2-(e_1-e_2)(e_1-e_3)\right)^2.
\end{equation}
当然，若 $w=e_2$ 或 $w=e_3$ 也有类似的公式。(3.16) 左端是完全平方这一事实意味着，在 $w=e_1$ 处，曲线 $\boldsymbol{F}_w$ 分裂为两个分支 $\boldsymbol{F}_{e_1}^{\pm}$ 的并，它们由
\begin{equation}\tag{3.17}
\rho=\pm\frac{1}{2}\left((u-e_1)^2-(e_1-e_2)(e_1-e_3)\right)
\end{equation}
定义。

曲线 $\boldsymbol{F}_{e_1}^{\pm}$ 各自亏格为 $0$。它们在由
\begin{equation}\tag{3.18}
\rho=0,\qquad u=e_1\pm\sqrt{(e_1-e_2)(e_1-e_3)}
\end{equation}
给出的两个点处相交。
% [Tutor Note: 推导补足。判别式之所以是 $\prod(e_i-e_j)^2\prod(w-e_i)^2$，只需两点：其一，$g(u)$ 是 $u$ 的四次多项式而系数关于 $w$ 线性，故其判别式是关于 $w$ 的 $6$ 次多项式；其二，$w=e_i$ 时 (3.16) 表明 $g$ 成为完全平方，即 $g$ 的两个根在 $u=e_1\pm\sqrt{(e_1-e_2)(e_1-e_3)}$ 处重合（这两点正是 (3.18) 的二重点），而「两根重合」使判别式（即诸根差平方之积）出现二重零点。三个 $e_i$ 各贡献一个二重零点，合计 $6$ 次，已与次数相符，故判别式恰为所给形式。 来源：原书 p.25 上下文（A）。]

一对在两处二重点相交的亏格 $0$ 曲线给出算术亏格 $1$ 的曲线。亏格 $1$ 的光滑曲线可以退化到这样的奇异曲线。这就是我们例子中 Hitchin 纤维化在三个纤维 $w=e_1,e_2,e_3$ 处的行为。这些奇异纤维显示在第 18 页的左图中。存在

% ==== 原文 PDF 第 26 页（印刷页 26）====

两个二重点（当 $w$ 等于某个 $e_i$ 时）正是 $g$ 的判别式在这些 $w$ 值处有二重零点的原因。

尽管 Hitchin 纤维化的纤维 $\boldsymbol{F}_w$ 在 $w=e_i$ 时奇异，$C$ 上稳定分歧 $SL_2$ Higgs 丛的模空间 $\mathcal{M}_H$ 实际上是光滑的。事实上，由于在 $\rho=0,w=e_i$ 处 $f\ne 0$，以 $\rho^2+(2u+w)f(u)-f'(u)^2/4$ 的零点刻画 $\mathcal{M}_H$ 的这个多项式在这些点处有非零微分。所以尽管 Hitchin 纤维化的纤维奇异，$\mathcal{M}_H(C;SL_2)$ 是光滑的。

我们所发现的正是第 2.2 节所预告的行为。Hitchin 纤维化的某些纤维分裂为两个分支 $\boldsymbol{F}^{\pm}$ 的并，从而导致 A-膜的分解。在第 3.6 节中，我们解释这与对偶群的模空间的奇点有何关系。

\subsection{对称群}

$SL_2$ Higgs 丛 $(E,\varphi)$ 可以被一个 $2$ 阶线丛 $\mathcal{N}$ 扭转。更确切地说，这一操作是 $E\to E\otimes\mathcal{N}$，$\varphi\to\varphi$。对亏格 $1$ 的 $C$，$2$ 阶线丛构成的群是 $\mathcal{Q}=\mathbb{Z}_2\times\mathbb{Z}_2$，而这个群必须作用在分歧稳定 Higgs 丛的模空间 $\mathcal{M}_H$ 上。
% [Tutor Note: 推导补足。「群必须作用」的理由是：若 $\mathcal{N}$ 为 $2$ 阶线丛，$E\otimes\mathcal{N}\cong E$（因为 $\mathcal{N}^2\cong\mathcal{O}$ 且秩 $2$ 丛张量 $\mathcal{N}$ 后同构类不变），故扭转只改变 Higgs 场的基，而不改变 $(E,\varphi)$ 的同构类；但不同的 $\mathcal{N}$ 给出不同的自同构，于是 $\mathcal{Q}$ 作用在模空间上。对亏格 $1$，$2$ 阶线丛正是 $\mathrm{Pic}^0(C)$ 的 $2$-挠点，故 $\mathcal{Q}\cong C[2]\cong\mathbb{Z}_2\times\mathbb{Z}_2$，这一点在 Remark 3.1 中将与纤维上的平移作用联系起来。 来源：原书 p.26 上下文（A）。]

底丛 $E$ 是 $E=\mathcal{L}\oplus\mathcal{L}^{-1}$，其中在对换 $\mathcal{L}\leftrightarrow\mathcal{L}^{-1}$ 之下，$\mathcal{L}$ 由 $u=x_0$ 参数化。群 $\mathcal{Q}$ 通过半稳定丛模空间的全纯自同构作用在 $E$ 上。这样的自同构是 $u$-平面的分式线性变换。此外，$\mathcal{L}$ 为 $2$ 阶这一条件在 $\mathcal{Q}$ 的作用下不变。

若 $u=\infty$ 则 $\mathcal{L}$ 平凡；若 $u=e_1,e_2$ 或 $e_3$ 则 $\mathcal{L}$ 是非平凡的 $2$ 阶线丛。因此用 $2$ 阶线丛扭转把 $u=\infty$ 与某个值 $u=e_i$ 互换，同时也把另外两个 $e$ 互换。例如，存在元素 $T_1\in\mathcal{Q}$，它把 $\infty$ 与 $e_1$ 互换，并且也把 $e_2$ 与 $e_3$ 互换。它作用为
\begin{equation}\tag{3.19}
\begin{aligned}
u &\to \frac{e_1u+e_2e_3-e_1e_3-e_1e_2}{u-e_1}\\
w &\to w\\
\rho &\to -\frac{(e_1-e_2)(e_1-e_3)}{(u-e_1)^2}\,\rho.
\end{aligned}
\end{equation}
$w$ 不变这一事实反映了：$\varphi$（从而其特征多项式）在用 $2$ 阶线丛扭转下不变。$\rho$ 变换中的符号并不显然（定义 $\mathcal{M}_H$ 的方程 (3.15) 在 $\rho\to-\rho$ 下不变），它可以由 $\mathcal{M}_H$ 的全纯二形式 $\Omega_I$ 是 $\mathcal{Q}$-不变这一事实确定。$\mathcal{Q}$ 还有由 $e_1,e_2,e_3$ 的循环置换得到的元素 $T_2$ 与 $T_3$。这些就是 $\mathcal{Q}$ 的三个非平凡元素。
% [Tutor Note: 推导补足。(3.19) 中 $u$ 的变换是使 $\{e_1,\infty\}$ 与 $\{e_2,e_3\}$ 各自作为集合不变的唯一分式线性对合（由交比确定）。验证 $\rho$ 的变换需要一步初等计算：记 $u'=\frac{e_1u+e_2e_3-e_1e_3-e_1e_2}{u-e_1}$，则 $u'-e_1=\frac{(e_1-e_2)(e_1-e_3)}{u-e_1}$，于是 $(u'-e_1)^2-(e_1-e_2)(e_1-e_3)=-\frac{(e_1-e_2)(e_1-e_3)}{(u-e_1)^2}\left[(u-e_1)^2-(e_1-e_2)(e_1-e_3)\right]$，代入 (3.17) 即见 $\rho$ 的变换确实把 $\boldsymbol{F}_{e_1}^{+}$ 与 $\boldsymbol{F}_{e_1}^{-}$ 各自映到自身（符号与 (3.19) 一致）。 来源：原书 p.26 上下文（A）。]

由于 $w$ 是 $\mathcal{Q}$-不变的，群 $\mathcal{Q}$ 作用在 Hitchin 纤维化的每个纤维 $\boldsymbol{F}_w$ 上。让我们描述它在奇异纤维（比如 $w=e_1$ 处）上的作用。用上述公式作一简短计算表明，$T_1$ 把每个分支 $\boldsymbol{F}_{e_1}^{\pm}$ 映到自身，并固定式 (3.18) 的两个二重点。曲线 $\boldsymbol{F}_{e_1}^{\pm}$ 亏格为零，而 $T_1$ 在每条这样的亏格零曲线上作用为具有两个不动点的分式线性变换，这两个不动点即二重点。（关于 $\mathbb{C}\mathbb{P}^1$ 上具有两个不动点的对合，可考虑复 $z$-平面的变换 $z\to -z$，其不动点为 $0$ 与 $\infty$。）另一方面，$T_2$ 与 $T_3$ 交换两个分支 $\boldsymbol{F}_{e_1}^{\pm}$，并且也交换

% ==== 原文 PDF 第 27 页（印刷页 27）====

两个二重点。既然如此，$T_2$ 与 $T_3$ 自由地作用在奇异曲线 $\boldsymbol{F}_{e_1}$ 上，从而也自由地作用在 Hitchin 纤维化的所有邻近纤维上。

\bigskip
\noindent\textit{注 3.1.}\ $\mathcal{Q}$ 通过 $2$ 阶点群的平移作用在 Hitchin 纤维化的每个纤维 $\boldsymbol{F}_w$ 上。这在第 5.3 节所述的谱曲线构造中是清楚的。因此当 $\boldsymbol{F}_w$ 光滑时，$\mathcal{Q}$ 自由地作用。当 $\boldsymbol{F}_w$ 退化为两个分支 $\boldsymbol{F}_w^{\pm}$ 的并时，它有一个对应于在二重点处坍缩的单圈的「短」方向，以及一个互补的「长」方向。$T_2$ 与 $T_3$ 是长方向上的平移；它们自由地作用并交换两个二重点与两个分支。$T_1$ 是短方向上的平移，把每个分支映到自身，并以二重点为不动点。
